\documentclass[12pt]{article}
\usepackage[utf8]{inputenc}
\usepackage[T1]{fontenc}
\usepackage[margin=1in]{geometry}
\usepackage{setspace}
\usepackage{xcolor}
\usepackage[colorlinks=true, urlcolor=blue!60!black, linkcolor=blue!60!black]{hyperref}
\setlength{\parindent}{0em}
\setlength{\parskip}{0.6em}
\renewcommand{\baselinestretch}{1.0}
\begin{document}
\begin{center}
{\LARGE \textit{Marbury v. Madison}, 5 U.S. (1 Cranch) 137 (1803)\par}
\vspace{0.5em}
{\large Case Brief --- Two-Page Edition\par}
\end{center}

\section{Memory Jogger}
Federal courts may refuse to enforce acts of Congress that conflict with the Constitution, and Congress cannot enlarge the Supreme Court's original jurisdiction (\href{https://supreme.justia.com/cases/federal/us/5/137/}{Justia}).

\section{Detailed Case Facts}
After losing the 1800 election, John Adams appointed forty-two District of Columbia justices of the peace, including William Marbury, to five-year terms. The Senate confirmed them, Adams signed the commissions, and Secretary of State John Marshall sealed them, but Marbury's was undelivered when Thomas Jefferson took office. Jefferson told Secretary of State James Madison to withhold it.

\section{Procedural History}
Marbury petitioned the Supreme Court directly for a writ of mandamus under Section 13 of the Judiciary Act of 1789. After Congress cancelled the Court's 1802 terms, the case was argued in February 1803, with no one appearing for Madison.

\section{Judicial Votes}
4--0. Marshall wrote for Paterson, Chase, and Washington; Cushing and Moore did not participate. There were no separate opinions.

\section{Holding}
The petition was dismissed for lack of original jurisdiction. Article III fixes the Court's original jurisdiction, so Section 13 was void insofar as it enlarged it, and courts must refuse to apply a statute repugnant to the Constitution.

\section{Analysis of Opinions}
\textbf{Majority:} Marbury's appointment was complete when the President signed the commission, so his right had vested; where there is a legal right there is a legal remedy, at least for ministerial duties, though political acts are ``only politically examinable.'' But Article III's list of original-jurisdiction cases is exclusive, or it would be ``mere surplusage.'' A written constitution is ``the fundamental and paramount law of the nation''; the judicial oath and the Supremacy Clause confirm that courts must prefer the Constitution to a conflicting statute.

\section{Examples: Future Applications}
\textbf{Same-Side:} a statute giving the Supreme Court original jurisdiction over voters' suits against the Attorney General is void. \textbf{Opposite-Side:} mandamus to recognize a foreign republic is refused as a political question. \textbf{Fence-Sitter:} a President who obeys a judgment but enforces a similar law against non-parties raises the judicial-supremacy debate.

\section{Critique}
William W. Van Alstyne argued that Marshall should have recused himself, reached the merits improperly, and begged the question of who interprets the Constitution finally. Others, following Alexander Bickel, stress the countermajoritarian difficulty of judicial review.

\section{Key Quotations}
\begin{quote}
``It is emphatically the province and duty of the Judicial Department to say what the law is. Those who apply the rule to particular cases must, of necessity, expound and interpret that rule. If two laws conflict with each other, the Courts must decide on the operation of each.''
\end{quote}
\noindent --- Chief Justice Marshall, for the Court, 5 U.S. (1 Cranch) at 177.

\end{document}
